
\ifthenelse{\isundefined{\jfpda}}{
\begin{figure*}[t]
\begin{tabular}{cc}
\includegraphics[width=\columnwidth]{comparison-T50-2000iterations.png} &
\includegraphics[width=\columnwidth]{machines-exploits-T50.png}\\[-0.6cm]
(a) Attack quality comparison.  & (b) Runtime of 4AL. 
\end{tabular}
\vspace{-0.2cm}
\caption{Empirical results for 4AL compared to a global POMDP model.}
  \label{fig:experiments}
  \vspace{-0.5cm}
\end{figure*}
}{
  \begin{figure}[t]
    \begin{tabular}{cc}
      \includegraphics[width=0.5\linewidth]{comparison-T50-2000iterations.png} &
      \includegraphics[width=0.5\linewidth]{machines-exploits-T50.png}\\[-0.6cm]
      (a) Attack quality comparison.  & (b) Runtime of 4AL. 
    \end{tabular}
    \caption{Empirical results for 4AL compared to a global POMDP model.}
    \label{fig:experiments}
  \end{figure}
}


\section{Experiments}
\label{experiments}

We evaluated 4AL against the ``global'' POMDP model, encoding the
entire attack problem into a single POMDP. The experiments are run on
a machine with an Intel Core2 Duo CPU at 2.2 GHz and 3 GB of RAM.  The
4AL algorithm is implemented in Python.  To solve and evaluate the
POMDPs generated by Level 4, we use the APPL toolkit.\footnote{APPL
  0.93 at http://bigbird.comp.nus.edu.sg/pmwiki/farm/appl/}




\subsection{Test Scenario}
\label{experiments:test-suite}

Our test scenario is based on the network structure shown in
Figure~\ref{fig:scenario}. The attack begins from the Internet
(\start\ is the cloud in the top left corner). The network consists of
three areas---\emph{Exposed, Sensitive, User}---separated by
firewalls. Internally, each of Exposed and Sensitive is fully
connected (i.e., these areas are subnetworks), whereas User consists
of a tree of subnetworks separated by empty firewalls. Only two
machines are rewarded, one in Sensitive (reward 9000) and one in a
leaf subnetwork of User (reward 5000). The cost of port scans and
exploits is 10, the cost of OS detection is 50. We allow to scale the
number of machines $|M|$ by distributing, of every 40 machines, the
first one to Exposed, the second one to Sensitive, and the remaining
38 to User. The exploits are taken from Core Security's database. The
number of exploits $|E|$ is scaled by distributing these over 13
templates, and assigning to each machine $m$ one such template as
$I(m)$ (the known configuration at the time of the last pentest). The
initial belief $b_0(I,T)$, where $T$ is the time elapsed since the
last pentest, is then generated as outlined.

\begin{figure}[htb]
\vspace{-0.0cm}
\centering
\includegraphics[width=0.95\columnwidth]{ff-test-scenario3.png}
  \vspace{-0.3cm}
\caption{Network structure in our test suite.}
  \label{fig:scenario}
  \vspace{-0.6cm}
\end{figure}

The fixed parameters here (rewards, action costs, distribution of
machines over areas, number of templates) are estimated based on
practical experiences at Core Security. The network structure and
exploits are realistic, and are used for industrial testing in that
company. The main weakness of the scenario lies in the approximation
of software updates underlying $b_0(I,T)$. Altogether, the scenario is
still simplified, but is natural and does approach the complexity of
real-world penetration testing.

For lack of space, in what follows we scale only $|M|$ and $|E|$,
fixing $|T| = 50$. The latter is realistic but challenging: pentesting
is typically performed about once a month; smaller $T$ are easier to
solve as there is less uncertainty.


\subsection{Approximation Loss}
\label{experiments:approximation}

Figure~\ref{fig:experiments} (a) shows the relative loss of quality
when running 4AL instead of a global POMDP solution, for values of
$|E|$ and $|M|$ where the latter is feasible. %
We show
$\mathit{quality}(\mathit{global\mbox{-}POMDP})-\mathit{quality}(\mathit{4AL})$
in percent of $\mathit{quality}(\mathit{global\mbox{-}POMDP})$. Policy
quality here is estimated by running 2000 simulations. That
measurement incurs a variance, which is almost stronger than the very
small quality advantage of the global POMDP solution. The maximal loss
for any combination of $|E|$ and $|M|$ is $14.1\%$ (at $|E|=7$,
$|M|=6$), the average loss over all combinations is $1.96\%$. The
average loss grows monotonically over $|M|$, from $-1.14\%$ for
$|M|=1$ to $4.37\%$ for $|M|=6$. Over $|E|$, the behavior is less
regular; the maximum average loss, $5.4\%$, is obtained when fixing
$|E|=5$.



\subsection{Scaling Up}
\label{experiments:scaling}

Figure~\ref{fig:experiments} (b) shows the runtime of 4AL when scaling
up to much larger values of $|E|$ and $|M|$. The scaling behavior over
$|M|$ clearly reflects the fact that 4AL is polynomial in that
parameter, except for the size of biconnected components (which is $3$
here). Scaling $E$ yields more challenging single-machine POMDPs,
resulting in a sometimes steep growth of runtime. However, even with
$|M|$ and $|E|$ both around $100$, which is a realistic size in
practice, the runtime is always below $37$ seconds.

